\documentclass[aps,pra,10pt,twocolumn,superscriptaddress]{revtex4-2}

\usepackage{amsmath,amssymb,amsfonts,mathtools,bm}
\usepackage{braket}
\usepackage{graphicx}
\usepackage{booktabs,array,tabularx}
\usepackage{microtype}
\usepackage{xcolor}
\usepackage[colorlinks=true,linkcolor=blue,citecolor=blue,urlcolor=blue]{hyperref}
\usepackage[capitalize]{cleveref}
\hypersetup{pdftitle={Top-down quantum state preparation under gate noise},
  pdfauthor={Maximilian Froehlich, Tomasz Szoldra, Martin Eigel, Michael Hintermueller, Peter Schmelcher},
  pdfsubject={First manuscript draft},
  pdfkeywords={quantum state preparation, circuit pruning, gate noise}}

\crefname{equation}{Eq.}{Eqs.}
\Crefname{equation}{Eq.}{Eqs.}
\crefname{section}{Sec.}{Secs.}
\Crefname{section}{Sec.}{Secs.}
\crefname{subsection}{Sec.}{Secs.}
\Crefname{subsection}{Sec.}{Secs.}
\crefname{figure}{Fig.}{Figs.}
\Crefname{figure}{Fig.}{Figs.}
\crefname{table}{Table}{Tables}
\Crefname{table}{Table}{Tables}
\crefname{appendix}{Appendix}{Appendices}
\Crefname{appendix}{Appendix}{Appendices}

\begin{document}

\title{Top-down quantum state preparation under gate noise}
% Author list and order supplied by Maximilian Froehlich.
\author{Maximilian Fr\"ohlich}
\affiliation{Weierstrass Institute for Applied Analysis and Stochastics, Berlin, Germany}
\author{Tomasz Szoldra}
\affiliation{Institute for Quantum Physics, University of Hamburg, Hamburg, Germany}
\author{Martin Eigel}
\affiliation{Weierstrass Institute for Applied Analysis and Stochastics, Berlin, Germany}
\author{Michael Hinterm\"uller}
\affiliation{Weierstrass Institute for Applied Analysis and Stochastics, Berlin, Germany}
\author{Peter Schmelcher}
\affiliation{Institute for Quantum Physics, University of Hamburg, Hamburg, Germany}

\date{September 20, 2026}

\begin{abstract}
Approximate quantum state preparation requires balancing circuit expressivity,
gate noise, and classical optimization cost. We study a top-down workflow that
constructs a parameterized circuit without noise, removes gates according to
their exact effect on target fidelity, and refines the surviving parameters
using frozen Pauli-noise trajectories. Noiseless construction combines
sequential adjoint updates with limited-memory quasi-Newton refinement. The
noisy phase uses a weighted cross-trajectory update without assuming that it
is the sampled-loss gradient. We compare the workflow with
three publication-based simulator adaptations and two internal controls on
nine fixed six-qubit targets under a synthetic logical-gate noise model. It
attains the highest selected fidelity on five targets and lies on the
fidelity--CPU Pareto frontier on seven. Its mean fidelity is 92.260\%, compared
with 91.019\% for RobustState, at 11.17\% greater construction and training CPU
cost. The strongest trade-offs occur on the tested real-valued targets.
Complex Haar targets and a separate native-gate study give less favorable
comparisons. These results identify target- and noise-model-dependent
trade-offs within the tested optimization budgets.
\end{abstract}

\maketitle

\section{Introduction}
\label{sec:introduction}

Preparing a prescribed quantum state is a basic component of quantum
simulation and variational quantum algorithms. General synthesis procedures
provide systematic constructions, but their circuit requirements grow rapidly
with the size of an arbitrary target~\cite{Mottonen2005,Plesch2011}.
Approximate preparation instead searches a restricted circuit family for a
state sufficiently close to the target. Its practical value depends on both
the resources required to execute the circuit and the classical work needed
to find it~\cite{Cerezo2021}.

Gate noise adds a second reason to control circuit size. Additional
parameterized gates can enlarge the set of accessible states, but executing
them also changes the noisy channel. The relation is not determined by gate
count alone. Different parameterizations of the same ideal state can have
different noise sensitivity, and redundant parameters can sometimes improve
the result of noisy reoptimization~\cite{Fontana2021}. Circuit reduction
therefore needs to be assessed through the fidelity of the resulting noisy
state, together with its construction cost.

Several approaches address parts of this problem. Noise-aware circuit
learning searches circuit structure and continuous parameters using a noisy
objective~\cite{Cincio2021}. Feedback-based refinement combines the measured
or simulated noisy output with derivatives of an ideal circuit
model~\cite{RobustState2024}. Parameter-efficient training alternates
optimization with pruning and regrowth~\cite{Sim2021PECT}. More recent
state-preparation work uses top-down sequential pruning, including selection
based on the cost of candidate gate removals~\cite{Ashfaq2026}.
Compilation-aware pruning also weighs the loss from omitting a rotation
against the noise associated with its routing overhead~\cite{Escofet2026}.
These methods establish circuit structure, parameter refinement, and noise
placement as coupled design choices.

Adjoint methods provide an efficient way to evaluate parameter sensitivities
in classical circuit simulation~\cite{JonesGacon2020}. Gate-local updates
inspired by Krotov optimal control use forward states and backward costates
to update parameters along the circuit~\cite{Goerz2019,Froehlich2026HybridKrotov}.
For the appropriate noiseless batch rule, this construction yields the
ordinary loss gradient. A sequential update uses changing forward states
while retaining backward information from the earlier iterate. These
ingredients can support circuit construction, but their use does not by
itself determine which gates to remove or how a sampled noisy refinement
should be interpreted.

Here, we study a specified combination of noiseless circuit growth,
target-dependent pruning, and noisy parameter refinement. At each visited
depth, sequential adjoint updates and limited-memory quasi-Newton refinement
construct an ideal circuit. We then evaluate the exact change in ideal
target fidelity caused by deleting each gate, remove gates below a fixed
threshold, and refine the surviving parameters using frozen noise histories.
Independent validation selects a depth before exact evaluation of the noisy
channel. The contribution is the construction and evaluation of this
workflow. We distinguish exact deletion scores and exact weighted history
compression from the implemented noisy update, for which we do not assume
an empirical-gradient identity or a monotonicity guarantee.

We evaluate the workflow on nine fixed six-qubit targets against three
publication-based simulator adaptations and two internal controls. Under
the synthetic logical-gate noise model, it gives the highest selected
fidelity on five targets and is Pareto-efficient in fidelity versus complete
construction and training CPU on seven. The real-valued targets provide the
strongest trade-offs. The complex Haar targets are more costly for the
workflow, and a separate native-gate comparison changes some rankings.
The results concern one optimizer seed and finite method-specific budgets;
they do not establish a general ranking of state-preparation algorithms.

We first define the state-preparation problem and gate-noise model. We then
describe the construction and selection procedure, followed by the numerical
protocol and the fidelity, runtime, and circuit-resource comparisons.
The discussion addresses the resulting scope and unresolved questions.

\section{State preparation and gate noise}
\label{sec:problem}

The objective is to approximate a known normalized target
$\ket{t}\in\mathbb C^{2^n}$ starting from $\ket{0^n}$.
For a circuit with $L$ elementary gates, we use the ordering
\begin{equation}
 C(\boldsymbol\theta)=G_L(\theta_L)\cdots G_1(\theta_1).
 \label{eq:circuit}
\end{equation}
Each trainable elementary gate has one scalar parameter. Fixed gates can
also occur in the comparator circuits. The noiseless fidelity and loss are
\begin{equation}
 F_0(\boldsymbol\theta)=
 |\bra{t}C(\boldsymbol\theta)\ket{0^n}|^2,
 \qquad \mathcal L_0=1-F_0.
 \label{eq:ideal_fidelity}
\end{equation}
All target amplitudes are available to the simulator. This access assumption
also permits exact evaluation of gate-removal damage.

After each elementary gate, a Pauli channel acts on its support. The channel
is parameter independent. For window $i$, let $P_{i,a}$ denote its nonidentity Pauli
outcomes and let $\lambda_{i,a}\geq0$ be their rates. With
$\kappa_i=\sum_a\lambda_{i,a}$, define
\begin{equation}
 p_{i,0}=e^{-\kappa_i},\quad
 p_{i,a}=(1-e^{-\kappa_i})\frac{\lambda_{i,a}}{\kappa_i}.
 \label{eq:jump_law}
\end{equation}
\begin{equation}
 \mathcal E_i(\rho)=p_{i,0}\rho+
 \sum_a p_{i,a}P_{i,a}\rho P_{i,a}^{\dagger}.
 \label{eq:channel}
\end{equation}
For $\kappa_i=0$, the channel is the identity. Each window contains either
no jump or one listed jump, and outcomes in different windows are sampled
independently. This defines the complete channel used here. It is not the
full repeated-jump evolution generated by a Lindblad equation.

For the profile used in the main comparison, each acted-on site has $X$,
$Y$, and $Z$ rates of $10^{-4}$ each. A two-qubit gate additionally has
correlated $XX$ and $ZZ$ rates of $1.5\times10^{-3}$ each. Thus
$\kappa_i=3\times10^{-4}$ for a one-qubit window and
$\kappa_i=3.6\times10^{-3}$ for a two-qubit window.
These coefficients are simulation inputs, not average gate-error
probabilities or a current device calibration. The historical profile name
is IBM. Idle noise and readout noise are absent.

Starting from $\rho_0=\ket{0^n}\bra{0^n}$, exact propagation gives
\begin{equation}
 \rho_i=\mathcal E_i\!\left(G_i\rho_{i-1}G_i^{\dagger}\right),
 \qquad F=\bra{t}\rho_L\ket{t}.
 \label{eq:noisy_fidelity}
\end{equation}
In the logical-gate study, a noise window follows every original elementary
operation, including a rotation whose angle is zero. Deleting a gate removes
both the operation and its noise window. Parameter changes alone do not
remove windows. Logical $R_{XX}$, $R_{YY}$, and $R_{ZZ}$ rotations and logical
CX gates all receive a two-qubit window. This convention fixes a simulation
model; their physical decomposition costs can differ.

The separate native study instead places the same arity-dependent channel
after operations in a compiled $\{R_Z,\mathrm{SX},X,\mathrm{CX}\}$ schedule.
The schedule is frozen during refinement. The two studies therefore answer
different implementation questions. Neither is a measurement on a quantum
processor.

Training propagates dense state vectors along sampled histories where the
method requires trajectories. Exact feedback methods and final assessment
use density matrices. At six qubits these objects contain 64 amplitudes and
$64\times64$ entries, respectively. The random-MPS target label describes
target generation, not the simulation backend.

\section{Preparation by construction, pruning, and noisy refinement}
\label{sec:method}

We first construct a circuit without noise, remove gates whose individual deletion has little effect on the target fidelity, and refine the remaining parameters under the prescribed channel. Circuit depth determines both the available parameter space and the number of noise windows. We therefore apply this procedure at successive depths and select a candidate using independent noisy validation. Figure~\ref{fig:workflow} summarizes the complete procedure used for the reported baseline.

\begin{figure*}[t]
    \centering
    \includegraphics[width=\textwidth]{figures/workflow.pdf}
    \caption{\textbf{Construction and assessment of the preparation circuit.} At each visited depth, noiseless optimization returns its best parameter vector. This unpruned circuit supplies the parameters for further growth; a separate copy is pruned and refined with frozen noise histories. Noisy refinement returns its final iterate. Independent validation controls depth selection and stopping, after which exact density-matrix propagation assesses the locked circuit. All visited depths contribute to construction and training cost, including discarded candidates.}
    \label{fig:workflow}
\end{figure*}

\subsection{Noiseless circuit construction}
\label{sec:construction}

The ansatz acts on a line of qubits. An initial product of single-qubit Euler rotations is followed by nearest-neighbor blocks arranged in alternating bond layers. We use the state-preparation adjoint of the nine-parameter brickwall matrix-product-disentangler block implemented in YAQS. Writing
\begin{equation}
    U(\vartheta,\varphi,\lambda)
    =R_z(\varphi)R_y(\vartheta)R_z(\lambda),
\end{equation}
one block has the form
\begin{equation}
    V=(U_a^\dagger\otimes U_b^\dagger)
    R_{ZZ}(-\gamma_z)R_{YY}(-\gamma_y)R_{XX}(-\gamma_x).
    \label{eq:preparation_block}
\end{equation}
The two Euler rotations and three entangling angles are independent. The implementation emits the blocks in reverse disentangling-layer order. The depth index $d$ counts complete even/odd bond cycles, giving $2d$ bond layers and $d(n-1)$ blocks. Before pruning, the circuit contains $3n+9d(n-1)$ scalar parameters. Each primitive rotation has its own parameter and noise window; an entire block is not treated as a single noisy operation.

We increase $d$ one cycle at a time and retain the previously trained coordinates. New coordinates are sampled from a centered normal distribution with standard deviation $0.05$ at the first depth and $0.001$ thereafter. All depths receive the same construction budget: 100 online adjoint sweeps followed by 100 L-BFGS updates. There is no ideal-fidelity stopping threshold within this search.

The online phase uses the gate-local discrete adjoint construction of the hybrid Krotov optimizer~\cite{Froehlich2026HybridKrotov}, motivated by the forward/backward structure of quantum optimal control~\cite{Goerz2019}. Let
\begin{equation}
    A_i=(\partial_{\theta_i}G_i)G_i^\dagger,
\end{equation}
and let $\ket{\psi_i}$ denote the state immediately after gate $i$. For the ideal infidelity $\mathcal L_0=1-F_0$, the terminal costate is
\begin{equation}
    \ket{b_L}=-\langle t|\psi_L^{\mathrm{old}}\rangle\ket{t}.
\end{equation}
We propagate this costate backward through the old circuit and evaluate
\begin{equation}
    s_i^{(0)}=2\operatorname{Re}\langle b_i|A_i|\psi_i\rangle,
    \qquad
    \theta_i\leftarrow\theta_i-0.7s_i^{(0)}.
    \label{eq:ideal_signal}
\end{equation}
During an online sweep, backward states remain fixed while gates are updated in circuit order and each updated forward prefix is propagated onward. When all states instead use the same parameter vector, Eq.~\eqref{eq:ideal_signal} gives the exact noiseless loss gradient, as in reverse-mode circuit differentiation~\cite{JonesGacon2020}. The distinction is the parameter vector at which the forward and backward states are evaluated.

The batch phase applies L-BFGS~\cite{LiuNocedal1989} to this true gradient,
\begin{equation}
    p_k=-H_k\nabla\mathcal L_0(\theta_k),\qquad
    \theta_{k+1}=\theta_k+\alpha_kp_k,
    \label{eq:lbfgs_direction}
\end{equation}
with a bounded Armijo search and positive-curvature safeguards specified in Appendix~\ref{app:lbfgs}. After every complete sweep or batch update, we compare the ideal fidelity with the best observed value, including the initialization. At the end of construction, we restore the corresponding parameter vector; ties keep the earlier vector. This checkpoint rule protects the returned noiseless state from a worse final iterate. It does not require every online sweep to improve fidelity.

\subsection{Exact gate-removal scores}
\label{sec:pruning}

We next score the removal of each complete elementary gate from the restored noiseless circuit. The signed fidelity damage is
\begin{equation}
    \Delta_i=F_0(C)-F_0(C\setminus G_i).
    \label{eq:deletion}
\end{equation}
Our pruning rule, denoted B in the experiment archive, simultaneously removes all gates satisfying $\Delta_i<\varepsilon$, with $\varepsilon=0.01$. Surviving parameters retain their trained values. A negative score means that the individual deletion improves ideal fidelity. The score is computed from the target overlap, without a noisy objective or a derivative approximation.

A shared forward/backward calculation evaluates every individual removal in $O(L2^n)$ work and storage for dense states; Appendix~\ref{app:deletion} gives the contraction. These individual scores do not bound the damage from simultaneous deletions. We neither sum them into a fidelity guarantee nor repair the circuit between removals. Each deleted gate also loses its noise window. The resulting balance between ideal approximation error and noise exposure is assessed after refinement and validation.

\subsection{Frozen trajectories and noisy refinement}
\label{sec:refinement}

For each pruned circuit we sample $M=100$ complete histories from the gate-local law and keep them fixed throughout refinement. The histories are parameter independent because the Pauli probabilities depend on gate support rather than rotation angle. Reusing them provides common random numbers across parameter updates.

If $m_0$ sampled histories contain no jumps, we merge them into one representative with weight $m_0/M$. Every history containing a jump retains weight $1/M$, including histories whose jumps may cancel. With $K=M-m_0+\mathbf 1_{m_0>0}$ retained histories, the sampled fidelity is
\begin{equation}
    \widehat F(\theta)=\sum_{r=1}^{K}w_r
    |\langle t|\psi_r(\theta)\rangle|^2.
    \label{eq:weighted_fidelity}
\end{equation}
States are recomputed after parameter changes. This compression preserves the original empirical weights and both implemented update modes exactly, as shown in Appendix~\ref{app:compression}. A nonempty circuit whose complete bank contains no jumps is recorded as a coverage failure; the bank is neither replaced nor conditioned on observing a jump.

Let $\ket{\psi_{s,i}}$ be the forward state of history $s$ immediately after gate $i$, before its noise. Let $\ket{\chi_{r,i}}$ be the unscaled target propagated backward through the suffix of history $r$, including the adjoint of the noise following gate $i$, to that same point. The implemented noisy signal is
\begin{equation}
\begin{split}
    s_i^{(\times)}={}&-2\operatorname{Re}
    \sum_{r,s=1}^{K}w_rw_s
    \langle\chi_{r,i}|A_i|\psi_{s,i}\rangle\,
    \langle\psi_{s,i}|\chi_{r,i}\rangle.
\end{split}
    \label{eq:cross_signal}
\end{equation}
Both trajectory indices carry empirical weights. We use this cross-trajectory signal as an update rule; we do not identify it with the gradient of $1-\widehat F$. Consequently, the noiseless exact-gradient result does not supply a convergence or monotonicity guarantee for this phase.

Refinement performs 100 online sweeps followed by 100 simultaneous updates, each with $\theta_i\leftarrow\theta_i-0.2s_i^{(\times)}$. Online sweeps reuse old backward states while refreshing forward prefixes, as in construction. Simultaneous updates evaluate every signal at the same iterate. There is no learning-rate decay, L-BFGS step, or best-iterate restoration in this phase: the historical refinement procedure returns its final parameters.

\subsection{Depth selection and final assessment}
\label{sec:depth_selection}

Every visited depth produces a pruned, refined candidate, while the next construction starts from the unpruned noiseless checkpoint. Independent validation under the same channel determines whether to continue growing. A material improvement must exceed $0.001$ plus twice the combined standard error of the current candidate and the last material record. We stop after three unsuccessful extensions or at $d=16$; the precise rule is given in Appendix~\ref{app:selection}.

The returned circuit is the candidate with greatest validation fidelity over all valid visited depths, with earlier depth preferred on a tie. This winner can differ from the last material record. The archived baseline used 32,768 validation histories per candidate, and its selection remains fixed in the present comparison. The later 8,192-history comparison validation is diagnostic for this endpoint. Exact density-matrix evaluation is performed only after selection is locked and never replaces the selected circuit. Thus the reported fidelity integrates all modeled noise outcomes, while the depth decision retains its original sampled-validation provenance.

\section{Numerical protocol}
\label{sec:protocol}

The numerical study asks whether the complete construction and refinement
workflow provides a useful combination of noisy fidelity, classical cost,
and circuit resources. We compare the frozen top-down method with three
publication-based procedures and two internal controls. Candidate grids and
optimizer budgets are fixed for each procedure, but are not equal across
procedures. The comparison therefore measures the implementations at their
declared budgets.

\subsection{Targets and comparison methods}

The target set contains nine six-qubit states. One encodes the square root
of a Gaussian probability density with mean $0.5$ and standard deviation
$0.1$ on an endpoint-excluded quantics grid. Three are ground states of the
open transverse-field Ising chain
\begin{equation}
 H=-\sum_{j=0}^{4}Z_jZ_{j+1}-h\sum_{j=0}^{5}X_j,
 \label{eq:tfim}
\end{equation}
with $h=0.5$, $1$, and $1.5$. We call these ferromagnetic, critical, and
paramagnetic, respectively; critical refers to the thermodynamic critical
ratio in a finite chain. Three further targets are independent complex Haar
states. The remaining two are normalized real random MPS with maximum bond
dimensions two and three. The fixed amplitudes, phase convention, and
generation seeds define the instances. All comparisons use optimizer seed
300. The target fixture was used in developing the top-down method and is
not an untouched algorithm-development holdout.

\Cref{tab:protocol} summarizes the candidate and training policies.
NACL searches structure and angles with a noisy objective, using OrthoMADS
for continuous optimization~\cite{Cincio2021}. Its checkpoints at steps 0,
10, and 50 belong to one search chain. RobustState retains the published
combination of a noisy-output residual and a noiseless circuit
Jacobian~\cite{RobustState2024}. We provide its feedback through exact
density matrices, without finite-shot measurement or hardware calibration.
The Fontana adaptation uses real $R_Y$/CX layers and noiseless BFGS followed
by noisy BFGS~\cite{Fontana2021}. Its real-state restriction excludes the
three complex Haar targets. The Gaussian, Ising, and two MPS fixtures are
all supported.

\input{tables/protocol}

These three comparators are local implementations of published procedures
with documented adaptations. They are not the original authors' software
or reproductions of their device experiments. The fixed VQSP control
optimizes independently initialized circuits at two depths. The unpruned
layerwise control grows a circuit through four depths. These controls
evaluate alternative complete workflows; neither is an operator-pool
ADAPT-VQE implementation. AQCE and BAA-LRSP are outside the chosen
noise-aware training scope. IQSP was excluded from this six-qubit comparison
because the then-current local reconstruction failed during initialization.
That failure provides no evidence about the published method's performance.
The later twelve-qubit IQSP diagnostics discussed in \cref{sec:discussion}
are separate from this ranked comparison.

\subsection{Selection and exact assessment}

Training, candidate selection, and final assessment use separate roles for
the noise model. The top-down and internal-control trajectory banks contain
100 unconditional histories and remain fixed during each refinement.
Identical no-jump paths retain their empirical probability weights.
Nonempty candidates whose sampled banks contain no jumps fail the coverage
check; sampling is not conditioned or repeated to force a jump.

The top-down endpoints were previously selected by the depth-search rule
in \cref{sec:method}, using 32,768 independent validation histories per
depth. Their archived selections remain fixed. The present comparison uses
an independent 8,192-history validation bank for each assessed candidate;
for a frozen top-down endpoint this is a diagnostic. Other methods select
the eligible candidate with greatest validation fidelity within the
declared grid. Ties are resolved by compiled CX count, CX depth, and
candidate identifier. The selection is recorded before exact assessment.

Final fidelities are obtained by exact density-matrix propagation of the
same logical channel with the selected parameters. This integrates all
noise outcomes and introduces no trajectory sampling error in the
reported endpoint score. It does not remove the randomness of training
or validation-based selection. Exact scores never replace a selected
circuit, even if another assessed candidate has a slightly higher exact
score. One optimizer seed per target does not support a seed-level
significance claim.

\subsection{Computational cost and implementation checks}

We charge each method for all visited construction and training candidates,
including discarded structural proposals, line-search work, and historical
computation reused from archives. Validation and exact final assessment are
recorded separately. Reusing a trained circuit avoids new computation but
does not make its construction free. The full comparison contains 51
supported method--target groups and 120 assessed circuit records, including
nine frozen top-down endpoints. That count does not include every
historically visited top-down depth. Their costs remain charged through the
archived search records.

A selected endpoint is Pareto-efficient in fidelity and CPU if no other
method on the same target has at least its fidelity and no greater
construction and training CPU, with one strict improvement. This definition
compares complete searches at their fixed budgets. Logical two-qubit count
and depth are reported separately from compiled CX resources. Logical
resource counts alone do not establish equal physical execution cost across
gate alphabets.

The computations use Python 3.11.9, NumPy 2.4.6, SciPy
1.17.1~\cite{Virtanen2020}, Qiskit 2.4.1~\cite{Qiskit2024}, and PyNomadBBO
4.6.0 for NACL. The all-target extension uses three independent target
processes with one numerical-library thread per worker. CPU is summed
over the charged work, rather than inferred from parallel wall time.
The archived manifests do not identify the host processor model, which
limits portability of the timing measurements. These are implementation
timings at six qubits, not hardware execution times or scaling measurements.

Independent checks reconcile all candidate records, selected endpoints,
cost totals, and Pareto sets. Stored source fingerprints and selection
locks identify the evaluated configuration. Direct channel propagation
also checks the final density matrices. These checks establish consistency
of the recorded calculation; they do not certify optimizer convergence or
broader benchmark generalization.

\section{Numerical results}
\label{sec:results}

We first compare the exact noisy fidelity of the validation-selected
circuits across all targets. We then include the classical and circuit
resources required to obtain those endpoints. The native-gate pilot tests
how much of the resulting interpretation survives a different placement
of the noise windows.

\subsection{Fidelity across target families}

\Cref{fig:fidelity} shows the complete logical-noise comparison. The
top-down workflow attains the highest selected-endpoint fidelity on five
of nine targets. Its mean fidelity is 92.260\%, compared with 91.019\% for
RobustState. The paired mean difference is 1.241 percentage points, with
five wins and four losses. These are descriptive comparisons of the
selected circuits, without an estimate of optimizer-seed variability.

\begin{figure*}[t]
 \centering
 \includegraphics[width=\textwidth]{figures/fidelity.pdf}
 \caption{\textbf{Selected fidelity under logical-gate noise.}
 (a) Exact density-matrix fidelity for the nine fixed six-qubit targets,
 optimizer seed 300, and the synthetic Pauli profile in
 \cref{eq:jump_law}. Each value refers to a circuit selected by independent
 validation, with archived top-down selections retained. Fontana excludes
 the three complex Haar targets. Bold entries mark the highest fidelity
 for each target. (b) Paired top-down minus RobustState differences in
 percentage points; blue and orange denote positive and negative values.
 The top-down workflow has the highest fidelity on five targets.
 Exact endpoint scores have no trajectory sampling error, while
 optimizer-seed variability remains unmeasured.}
 \label{fig:fidelity}
\end{figure*}

The six real-valued targets give the strongest results for the top-down
method. Their mean fidelity is 95.885\%, compared with 94.153\% for
RobustState. The top-down method wins on Gaussian, ferromagnetic and
critical Ising, and bond-two MPS. RobustState is better on paramagnetic
Ising and bond-three MPS by 0.237 and 0.249 percentage points, respectively.
This division describes the tested fixture. It does not isolate real
amplitudes, entanglement, or another target property as the cause.

The complex Haar targets expose a different regime. RobustState is more
accurate on Haar 1 and Haar 3, with gaps of 1.132 and 0.551 percentage
points. On Haar 2, the top-down method instead gains 2.463 percentage
points. Thus even these three targets do not produce a uniform ranking.
Their top-down fidelities remain between 84.811\% and 85.250\%, below the
95.272--96.472\% range of the six real targets.

Fontana is supported on the six real targets only. On that common subset,
the top-down workflow wins in all six cases, with a mean advantage of
3.875 percentage points. Fontana's six-target mean must be interpreted on
this subset, rather than compared with nine-target means. NACL has lower
fidelity on all nine targets at the declared 50-step structural-search
budget. The result characterizes that finite search and does not establish
the performance of a converged NACL implementation.

The top-down workflow also exceeds both internal controls in fidelity on
every target. The mean improvements are 4.611 percentage points over
fixed VQSP and 2.929 over the unpruned layerwise control. The controls
differ in candidate grids and optimization policies, so these results
support the complete preparation workflow without assigning the gain
to pruning or noisy refinement alone.

\subsection{Classical cost and circuit resources}

\Cref{fig:pareto} combines selected fidelity with full construction and
training CPU. The top-down workflow is Pareto-efficient on seven of nine
targets, both against the external adaptations alone and when the internal
controls are included. Across the complete fixture it uses 285.17 CPU
seconds, compared with 256.52 for RobustState, an increase of 11.17\%.
The modest aggregate fidelity advantage therefore has an aggregate
classical cost.

\begin{figure*}[t]
 \centering
 \includegraphics[width=\textwidth]{figures/pareto.pdf}
 \caption{\textbf{Fidelity versus complete construction and training cost.}
 Each panel compares the selected endpoints for one target under the
 logical-gate noise model. CPU includes every visited candidate and the
 historical work of reused searches; it excludes validation and final
 assessment. The horizontal axes are logarithmic. The marked Pareto
 frontiers use all supported methods, including controls. Connections
 identify the discrete frontier and do not represent an interpolated
 performance curve. The top-down workflow is nondominated on seven
 targets; RobustState dominates it on Haar 1 and Haar 3.}
 \label{fig:pareto}
\end{figure*}

The cost separates sharply by target subset. On the six real targets,
top-down construction and training take 103.81 CPU seconds, compared with
170.92 for RobustState, a reduction of 39.27\%. Top-down is faster on each
of these targets. The two small fidelity losses to RobustState therefore
remain fidelity--runtime trade-offs. Fontana is substantially cheaper
still, taking 31.37 CPU seconds on the same six targets. The top-down
method's paired fidelity advantage over Fontana costs 3.31 times as much
classical computation.

RobustState dominates top-down on Haar 1 and Haar 3 in both fidelity and
CPU. It also uses fewer logical two-qubit gates and less two-qubit depth.
On Haar 2, the top-down fidelity gain costs 53.16 rather than 29.35 CPU
seconds and requires 38 rather than 20 logical two-qubit gates. Higher
fidelity on this target is therefore obtained with more classical work
and a larger logical circuit.

The fixed VQSP control is cheaper overall, at 115.88 CPU seconds, while
the layerwise control takes 354.24 seconds. The latter has both lower
mean fidelity and greater aggregate CPU than top-down. Aggregate values
do not replace the per-target Pareto comparisons. Full fidelity and
resource tables are included in the appendices.

These costs charge the complete method histories. The extension campaign
itself took 374.21 seconds of wall time and performed 966.96 CPU seconds
of fresh construction and training. Including historical work, the
charged total across all methods is 1,753.68 CPU seconds. The campaign
wall time is consequently not a measure of any single method's cost.

\subsection{Native-gate noise}

The logical comparison assigns one arity-dependent window to every
original elementary gate. To examine the effect of native implementation,
we also report the matched three-target pilot in \cref{tab:native}.
Its circuits are trained, selected, and assessed with noise after a frozen
$R_Z$/SX/$X$/CX schedule. The top-down method reuses the same noiseless
construction and pruning B checkpoints, followed by native refinement.
All historical construction work is charged.

\input{tables/native_results}

The native comparison is less favorable to top-down. On Gaussian,
Fontana reaches 94.144\% fidelity at 19.66 CPU seconds, compared with
92.891\% at 26.17 seconds for top-down. On Haar 1, RobustState reaches
82.090\% at 52.24 seconds, compared with 70.220\% at 79.84 seconds.
Both competitors also use fewer CX gates on the corresponding target.
These are direct counterexamples to a uniform native advantage.

Critical Ising retains a fidelity--runtime trade-off. Top-down reaches
90.062\%, compared with 90.334\% for RobustState, while taking 33.46
rather than 52.29 CPU seconds. Its 36.0\% CPU saving accompanies a
0.272-percentage-point fidelity loss and 26 rather than 10 CX gates.
Fontana is faster than either at 9.65 seconds, with 88.005\% fidelity.

After native selection was locked, a separate diagnostic evaluated the
earlier logically trained parameters on the same frozen native schedules
at the selected depths. Relative to those parameters, native refinement
changed fidelity by $-0.021$, $-0.026$, and $+0.020$ percentage points for
Gaussian, critical Ising, and Haar 1. The small changes do not show a
substantial benefit from changing the training noise location in this
pilot. They also should not be confused with the larger differences from
earlier numerically compiled deployment circuits, which involve schedule
changes and, for Haar 1, a different selected depth.

Native refinement nevertheless improves upon the pruned starting circuits
by 0.974, 4.446, and 3.842 percentage points, respectively. The absence of
a further gain over logical training on the matched schedules is a
different observation from the gain obtained by refining the pruned input.
The pilot uses one seed and three targets, with 32,768-path top-down
validation and the competitors' retained 8,192-path selections. It limits
the interpretation of the logical results without establishing a general
native-gate ranking.

\section{Discussion}
\label{sec:discussion}

We studied a top-down preparation workflow that separates ideal circuit
construction, gate removal, and noisy refinement. Exact single-gate
deletion scores identify the effect of each removal on the current ideal
target fidelity. Frozen Pauli histories make the sampled refinement
repeatable, and empirical weighting preserves the contributions of
duplicate no-jump paths. Independent validation selects the circuit
before exact evaluation of the same channel. These operations give a
specified and auditable preparation procedure.

The distinctions between its stages matter. Exact deletion damage for
one gate does not bound the combined damage from simultaneous removals.
Restoring the best noiseless parameters protects the returned ideal
checkpoint, while the noisy phase retains its final iterate. Its
cross-trajectory signal is not identified with the sampled-loss gradient.
The noiseless adjoint interpretation and quasi-Newton safeguards therefore
do not provide a convergence or monotonicity guarantee for the complete
noisy workflow. The numerical results evaluate the implemented procedure
with these limitations intact.

Under logical-gate noise, the method combines the highest selected fidelity
on five of nine targets with Pareto efficiency in fidelity and classical
cost on seven. The six real targets give the strongest trade-offs against
RobustState. Fontana remains an inexpensive alternative on that subset.
Complex Haar targets require more top-down construction work, and
RobustState dominates two of the three instances. The internal controls
support the usefulness of the full workflow at the tested settings, but
their differing policies prevent attributing the observed gains to a
single component.

Noise placement is a substantive part of this conclusion. Logical
two-qubit primitives do not have equal physical decomposition costs.
The native pilot gives a weaker top-down comparison on Gaussian and
Haar 1, while retaining a classical-cost trade-off on critical Ising.
The same-schedule parameter diagnostic also gives little benefit from
switching the refinement to native noise at the selected endpoints.
Thus the favorable logical results should be interpreted together with
the native results rather than extrapolated to a common physical gate
implementation.

RobustState's refinement benefit also depends on the noise mechanism.
Its authors attribute the hardware gains primarily to coherent-error
correction and report only marginal fidelity improvement in their
incoherent-noise simulator~\cite[Sec.~IV.C]{RobustState2024}.
Our Pauli-noise benchmark excludes systematic coherent gate errors and
therefore probes a related regime in which additional benefit from
parameter refinement may be limited. This qualifies the expected
refinement gain, rather than the achievable preparation fidelity:
RobustState remains competitive on our complex Haar targets.
The present comparison does not assess its reported advantage in
correcting coherent hardware errors.

Separate twelve-qubit diagnostics indicate a substantial cost for noisy
IQSP in our benchmark implementation. An early depth-four Gaussian
checkpoint required $505\,\mathrm{s}$ for one full fidelity-gradient
evaluation with our default compression on one Apple M1 Pro thread; peak
physical footprint was $4.76\,\mathrm{GiB}$.
% Provenance and compression settings: data/q12_iqsp_resource_diagnostic.json.
The published H2 demonstration uses noiseless one-qubit gates and
angle-dependent, factorized single-qubit depolarization after native $ZZ$
rotations~\cite{Ballarin2025}. Our benchmark additionally includes joint
two-qubit Pauli errors, requiring joint noise contractions and compression
in the locally purified implementation. Our trajectory refinement instead
propagates sampled pure states without purification indices, at the cost
of sampling error. This suggests a potential computational advantage for
parameter-independent Pauli noise. These checkpoint diagnostics neither
isolate noise structure as the cause of the observed cost nor establish
superior end-to-end performance or scaling. IQSP itself is not restricted
to the noise model or circuits used in its demonstration.

The main comparison is limited to a development-used fixture, six qubits, one
optimizer seed, one synthetic Pauli profile in the all-target comparison,
and finite, unequal search budgets. The exact-density feedback supplied
to several comparators is a simulator adaptation. It does not measure
their hardware sampling costs or their response to coherent device
errors. The dense implementation also does not establish the scaling of
a tensor-network realization. In particular, the retained-history cross
update contains contractions over pairs of trajectories; duplicate
compression alone does not imply linear scaling with the raw sample
count.

Further evaluation should vary optimizer seeds and use new target
instances, extend the matched native study, and test budget sensitivity.
Matched component ablations would be needed to separate the roles of
pruning, construction refinement, and the noisy update. Mathematical
analysis of the cross-trajectory signal is a separate question from
implementation correctness. Within the measured setting, top-down
pruning and refinement provide a useful preparation option whose value
depends on the target, circuit representation, and available classical
computation.


\section*{Code and data availability}
The implementation is developed in MQT YAQS~\cite{YAQS,Will2024MQT}.
The accompanying manuscript source package contains the data used for its
figures and tables, plotting scripts, and source checksums. The full training
histories remain in the experiment archives, including earlier archives
referenced by the final comparison. Those dependencies are listed in the
package's provenance notes. The manuscript package is sufficient to compile
the paper and regenerate its figures; reproducing the optimization runs also
requires the archived sources, histories, and pinned experiment environment.
No public deposit identifier for that complete experiment archive is assigned
in this draft.

\bibliography{references}

\appendix
\section{Shared evaluation of individual gate removals}
\label{app:deletion}

The exact deletion score requires the target overlap after omitting one gate, with every other gate and parameter fixed. Define the prefix state and suffix target
\begin{align}
    \ket{u_{i-1}}&=G_{i-1}\cdots G_1\ket{0^n},\\
    \ket{v_i}&=G_{i+1}^\dagger\cdots G_L^\dagger\ket{t}.
\end{align}
At the boundaries, $\ket{u_0}=\ket{0^n}$ and $\ket{v_L}=\ket{t}$. Their overlap gives
\begin{equation}
    F_0(C\setminus G_i)=|\langle v_i|u_{i-1}\rangle|^2.
    \label{eq:deletion_overlap}
\end{equation}
One forward pass stores the prefixes. A backward pass contracts each prefix with the corresponding suffix target before applying $G_i^\dagger$ to continue the recursion. Each gate has support on at most two sites, so its dense application and each overlap require $O(2^n)$ operations. Computing all $L$ scores therefore takes $O(L2^n)$ work with an $O(L2^n)$ prefix tape. No complete circuit is propagated separately for every deletion.

Equation~\eqref{eq:deletion_overlap} evaluates individual omissions from the same circuit. The simultaneous rule uses those scores once, removes the eligible gates, and compacts the surviving parameter indices. It does not equate the fidelity of the jointly pruned circuit with a sum of individual changes. This distinction matters because gates can compensate for one another, even when each individual removal has small signed damage.

\section{Preservation under no-jump compression}
\label{app:compression}

Let $h=1,\ldots,M$ index the original equally weighted histories, and partition them into retained groups $\mathcal H_r$. The sole nonsingleton group, when present, contains the histories with no noise operators at any gate. All its members have identical forward states and backward targets for any common parameter vector. Set
\begin{equation}
    w_r=\frac{|\mathcal H_r|}{M}.
\end{equation}
For a trajectory fidelity $f_h$, exact grouping gives
\begin{equation}
    \frac1M\sum_{h=1}^{M}f_h
    =\sum_r w_rf_r.
    \label{eq:compression_linear}
\end{equation}
The same reasoning applies to the pair contribution
\begin{equation}
    Q_{hh',i}=\langle\chi_{h,i}|A_i|\psi_{h',i}\rangle
    \langle\psi_{h',i}|\chi_{h,i}\rangle.
\end{equation}
Because each factor depends only on the corresponding retained history,
\begin{equation}
    \frac1{M^2}\sum_{h,h'=1}^{M}Q_{hh',i}
    =\sum_{r,s}w_rw_sQ_{rs,i}.
    \label{eq:compression_bilinear}
\end{equation}
Equation~\eqref{eq:compression_bilinear} explains why weighting only one index would change the cross update. Group members also remain identical during an online sweep: their noise maps coincide and they share every parameter update. Compression therefore preserves that sweep as well as simultaneous updates, up to floating-point arithmetic. It reduces repeated work but neither changes the sampled channel law nor establishes that the cross signal differentiates Eq.~\eqref{eq:weighted_fidelity}. The number of raw histories remains $M=100$, irrespective of $K$.

\section{Safeguards in noiseless L-BFGS refinement}
\label{app:lbfgs}

The L-BFGS history begins empty after the online phase at every depth. Let $g_k=\nabla\mathcal L_0(\theta_k)$, and form pairs from successive accepted states,
\begin{equation}
    s_k=\theta_{k+1}-\theta_k,\qquad y_k=g_{k+1}-g_k.
\end{equation}
A pair is retained only when both vectors and their norms are finite, both norms are nonzero, and
\begin{equation}
    s_k^Ty_k>10^{-10}\|s_k\|\|y_k\|.
\end{equation}
At most ten pairs enter the usual two-loop recursion~\cite{LiuNocedal1989}. Its initial inverse-Hessian scale is the latest $s_k^Ty_k/(y_k^Ty_k)$ clipped to $[10^{-8},10^8]$, or one for empty history. If the resulting direction or its slope is nonfinite, or $g_k^Tp_k\geq0$, the history is cleared and the direction becomes $-g_k$.

The initial trial step is $0.7$ with empty history and one otherwise. We accept the first finite trial satisfying
\begin{equation}
\begin{split}
    \mathcal L_0(\theta_k+\alpha p_k)
    \leq\mathcal L_0(\theta_k)
    +10^{-4}\alpha g_k^Tp_k.
\end{split}
    \label{eq:armijo}
\end{equation}
There are at most twelve halvings after the initial trial, hence at most thirteen trials per update. Failure retains the current vector and consumes one of the 100 batch-update slots. Trial evaluations use the current gradient without recalculating it at rejected points. New curvature information is obtained from the next batch gradient after an accepted step. Parameters are not wrapped, and there is no gradient-norm early exit. Best-iterate retention spans the initialization, the complete online sweeps, and the batch endpoints; rejected probes are not checkpoints.

\section{Validation, stopping, and the selected depth}
\label{app:selection}

For each valid refined baseline candidate, independently sample $M_v=32768$ unconditional validation histories. If $f_r$ are retained trajectory fidelities and $v_r$ their empirical multiplicity weights, the mean and estimated squared standard error are
\begin{align}
    \widehat F_d^{\mathrm{val}}&=\sum_rv_rf_r,\\
    \mathrm{SE}_d^2&=\frac{\sum_rv_r(f_r-\widehat F_d^{\mathrm{val}})^2}{M_v-1}.
    \label{eq:validation_se}
\end{align}
The denominator uses the raw sample count. Compression does not turn the retained representatives into equally weighted independent observations.

The first valid candidate establishes a material record $(F_*,\mathrm{SE}_*)$. A subsequent candidate replaces this record and resets patience only if
\begin{equation}
    \widehat F_d^{\mathrm{val}}>F_*+0.001
    +2\sqrt{\mathrm{SE}_d^2+\mathrm{SE}_*^2}.
    \label{eq:material_record}
\end{equation}
Otherwise patience increases by one; a coverage failure also counts as an unsuccessful extension. The search stops at patience three or depth sixteen. Equation~\eqref{eq:material_record} is a stopping heuristic with an uncertainty allowance, not a formal sequential significance test.

The selected depth is separately maintained as
\begin{equation}
    d^*=\underset{d\;\mathrm{valid\ and\ visited}}{\operatorname{argmax}}
    \widehat F_d^{\mathrm{val}},
    \label{eq:depth_winner}
\end{equation}
with the earlier depth retained on equality. Even an improvement too small to reset patience can change this winner. The archived selection is preserved when its trained circuit enters the present comparison. Other methods select within their declared grids using their independent 8,192-history validation and the fixed resource tie rule described in the experimental protocol. Exact final evaluation integrates the same channel for fixed parameters and reports no trajectory sampling error; it does not remove variability from initialization, training histories, or depth selection.

\section{Target conventions and complete endpoint data}
\label{app:targets}

The target fixture stores normalized amplitudes in little-endian basis
order: bit $j$ of the amplitude index is qubit $j$. The Gaussian coordinate
uses the quantics convention
\begin{equation}
\begin{aligned}
 x&=\sum_{j=0}^{5}s_j2^{-(j+1)},\\
 t(s_0,\ldots,s_5)&\propto
 \exp\!\left[-\frac{(x-0.5)^2}{4(0.1)^2}\right].
\end{aligned}
\end{equation}
Thus qubit zero has the largest coordinate weight even though it is the
least significant bit of the stored amplitude index. The amplitudes are
square roots of sampled Gaussian probabilities and are normalized over
the finite grid.

The Ising targets are the lowest eigenstates of Eq.~\eqref{eq:tfim}.
Each Haar target is obtained by drawing independent standard normal real
and imaginary parts and normalizing. Their generation seeds are 4001,
4002, and 4003. The MPS tensors contain independent standard normal real
entries, with bond sequences $(1,2,2,2,2,2,1)$ and $(1,3,3,3,3,3,1)$ and
seeds 5002 and 5003. Contracting and normalizing produces the dense target
vectors~\cite{Schollwock2011}. Every target's global phase is fixed so
that a largest-magnitude amplitude is positive real. The full fixture
also contains larger targets, which are not part of this study.

\Cref{tab:full-fidelity} provides the complete selected fidelities.
\Cref{tab:resources} records the corresponding resource comparison for
top-down and RobustState. The machine-readable endpoint data accompany
the manuscript and retain all methods' resources and CPU costs.

\input{tables/full_fidelity}
\input{tables/resources}

\section{Comparator adaptations and selection diagnostics}
\label{app:comparators}

NACL starts from an empty circuit and uses directed nearest-neighbor
CX gates, $R_Z$ rotations, and SX gates. Structure proposals include
one- and two-qubit identity insertion, gate removal, and adjacent
shuffling. Continuous angles are reoptimized using OrthoMADS, with a
maximum of 100 objective evaluations per inner solve. The implementation
uses one chain, an initial annealing temperature of $0.01$, a factor
$0.95$ per structural step, and simplification every five steps.
Proposal probabilities and these search conventions are documented
reconstruction choices. The step-zero candidate contains no gates,
and therefore has no noise windows or jump histories.

RobustState uses neighboring-pair CX blocks followed by local
$R_Z$/$R_X$/$R_Z$ rotations. For target projector $P=\ket{t}\bra{t}$,
its feedback loss and update signal are
\begin{align}
 L_{\rm R}&=\|\rho_{\rm feedback}-P\|_F,\\
 g_j^{\rm R}&=\operatorname{Tr}\!\left[
 \frac{\rho_{\rm feedback}-P}{L_{\rm R}}
 \frac{\partial\rho_{\rm ideal}}{\partial\theta_j}\right].
 \label{eq:robuststate_feedback}
\end{align}
The noisy residual is held fixed in this contraction with the ideal
Jacobian. It is not replaced by the true noisy-infidelity derivative.
At zero residual the signal is zero. Adam uses learning rate $0.005$,
coefficients $(0.9,0.999)$, and regularizer $10^{-8}$, with moments
retained across the ideal-to-noisy transition~\cite{Kingma2015}.
The final iterate is returned. Exact-density feedback substitutes for
the measurement procedure; the original coherent-error and hardware
sampling claims are not tested here~\cite{RobustState2024}.

The Fontana ansatz extends the published four-qubit real-state layout
to six qubits on a line~\cite{Fontana2021}. Each layer applies one
$R_Y$ to every qubit, followed by even and odd neighboring CX pairs.
Depth $d$ has $6d$ parameters and $5d$ CX gates. Each depth has an
independent uniform-angle initialization. Both optimization phases
use SciPy BFGS with analytic derivatives, a 200-iteration limit,
gradient tolerance $10^{-8}$, and a joint loss/gradient-call cap of
$1+20\times200$ per phase. The returned vector is the best accepted
iterate, including the initial vector; rejected line-search probes are
not eligible. Its noise placement is the common logical channel in
the all-target comparison and the common native channel in the pilot,
not the original publication's layerwise noise model.

The fixed VQSP control uses L-BFGS-B~\cite{Byrd1995} at depths two and
four, with separate initializations and weighted 100-history noisy
objectives. The layerwise control retains the unpruned ansatz, online
Krotov/L-BFGS construction with its best parent, and online/cross noisy
refinement through depths one to four. These
controls share the target fixture and channel with the other methods,
but do not constitute matched single-component ablations.

As a selection diagnostic, RobustState at critical Ising has a depth-four
candidate whose exact fidelity exceeds that of the selected depth-two
candidate by 0.056805 percentage points. The independent validation
fidelity instead favors depth two by 0.089086 percentage points.
The original selection is retained. This example illustrates why
exact assessment and candidate selection are kept separate.

\section{Archive provenance and reproducibility}
\label{app:provenance}

The all-target comparison reuses 17 complete method--target groups from
the three-target logical pilot and six top-down endpoints from the
noiseless L-BFGS construction archive. The remaining 28 groups contribute
74 newly trained candidate records. Reused circuits retain their
parameters and selections; CPU includes their historical search costs.

The manifest pins 183 inputs, including source files, the target fixture,
and the logical-suite protocol, and records package versions. Independent
audits authenticated 153 construction, selection, and final-assessment
bundles and reconciled all 120 assessed candidate records. Every nonempty
validation bank contained jumps, with a minimum of 38 among 8,192 raw
histories. Empty NACL circuits have no gate windows. These checks and the
targeted regression tests do not replace a multi-seed performance study.

The portable data include source hashes and reproduce the figures and
tables from the sealed logical and native reports. Replaying optimization
also requires the earlier construction archives, recursively referenced
histories, and pinned environment.

The historical pruning D, 400-history, and first-90\% studies are
excluded. The reported baseline uses pruning B, noiseless online
Krotov/L-BFGS at every visited depth, and 100-history noisy refinement.


\end{document}
